Table~\ref{tab:summary} sets the pilot against the five hypotheses of Section~\ref{sec:protocol}.
H1 to H4 were written down before any experiment was run.
H1 is given in the form to which the pilot led us to narrow it: originally it predicted that removing rule following would lower accuracy on \emph{every} exercise type, whereas \skill{induce}, which in TinyTextbook needs no table look-up, was still solved by \ETwoNoFollowInduceSolved{} of the \ETwoNoFollowSeeds{} seeds trained without \skill{follow}. The prediction is now restricted to skills that build on the one removed.
H5 was added during the pilot, after exploratory runs had already suggested its outcome, so it should be read as a question put to the full-scale study and not as a prediction that was tested fairly here.

\begin{table}[t]
\centering\small
\begin{tabular}{@{}lp{7.6cm}l@{}}
\toprule
 & \textbf{Evidence in the pilot} & \textbf{Status at this scale} \\
\midrule
H1 & Each skill left out of training stayed at chance in every seed; without \skill{follow}, \skill{compose} was solved in \ETwoPoolComposeWithoutFollowSolved{} of \ETwoPoolComposeWithoutFollowRuns{} runs, against \ETwoPoolComposeWithFollowSolved{} of \ETwoPoolComposeWithFollowRuns{} with it (E2). & Supported, after narrowing \\
H2 & Models that memorised their textbooks were at chance on new ones (E1). No direct test. & Not tested \\
H3 & At a fixed number of episodes, the number of distinct textbooks decided between memorising and reading (E1). Rule families were not varied. & Supported in part \\
H4 & With the same parameters, a second loop extended the chain length that was solved; a third made no visible difference. At equal depth, looped and unshared models showed no consistent difference (E3). & Supported in part \\
H5 & Separately trained skills did not combine: \skill{mixed} was solved in \ETwoPoolMixedAllSolved{} of \ETwoPoolMixedAllRuns{} runs, and chains longer than those trained stayed at the guessing level (E2, E3). & Contradicted \\
\bottomrule
\end{tabular}
\caption{The pilot against the hypotheses. ``At this scale'' means models of one to three hundred thousand parameters, one synthetic domain, one to four seeds.}
\label{tab:summary}
\end{table}
